 \section{Task Relaxation and Adaptation in Unknown Workspaces}
\label{ch03:sec:task_relaxation_adaptation}

The discrete planning framework of the previous chapter assumes that the
workspace abstraction is known before execution and remains fixed. This is
restrictive in partially known environments, where blocked passages, revised
semantic labels, or newly observed connectivity may change the feasibility of a
local temporal task. Planning must then determine not only a discrete plan, but
also how the task and model should be relaxed or updated as information arrives.
This section studies local temporal tasks under such incomplete workspace
knowledge. It first treats single-robot task relaxation, hard--soft
specifications, and online plan adaptation, and then extends the same ideas to a
multi-robot setting through collaborative knowledge transfer rather than a
centralized team specification.

%%%%%%%%%%%%%%%%%%%%%%%%%%%%%%%%%%%%%%%%%%%%%%%%%%%
\subsection{Infeasible tasks}\label{ch03:sec:infeasible}

The notation of Chapter~\ref{chap:preliminaries} is used throughout this
subsection. Consider a robot with local task specification~$\varphi$. By
Definition~\ref{ch02:def:feasible}, $\varphi$ may be infeasible over the
current abstraction~$\mathcal{T}_c$, either because the task is intrinsically
infeasible or because the model is incomplete. When the nominal synthesis
procedure of Section~\ref{ch02:sec:solution} fails, the objective becomes to
relax the specification and synthesize a plan that violates the original task as
little as possible.
A related approach in~\cite{kim2012approximate} constructs a relaxed
automaton~$\mathcal{A}_{\varphi}'$ that is close to
$\mathcal{A}_{\varphi}$ and feasible over~$\mathcal{T}_c$. The present
treatment uses a more discriminating objective: among accepting runs in the
relaxed product, select one that best preserves the original specification
under an explicit violation measure, and then infer the associated automaton
relaxation from that run.

\begin{problem}
Assume that the task specification is infeasible. How should the specification
be relaxed, and how can a discrete plan be synthesized so that it satisfies the
original task as much as possible?
\end{problem}

\subsubsection{Relaxed product automaton}
\label{ch03:subsec:relaxed-prod}

By Definition~\ref{ch02:def:feasible}, infeasibility of~$\varphi$ means that
the standard product automaton~$\mathcal{A}_p$ admits no accepting run. The
idea is therefore to relax the automaton-side constraints while keeping
the workspace unchanged. This leads to the relaxed product
B\"uchi automaton below.

\begin{definition}\label{ch03:def:relaxpro}
The relaxed product B\"uchi automaton is defined as:
\begin{equation}
\mathcal{A}_r \triangleq \mathcal{T}_c \times \mathcal{A}_{\varphi}
= (Q',\, 2^{AP},\, \delta',\, Q_0',\, \mathcal{F}',\, W_r),
\label{ch03:eq:relaxed-pba}
\end{equation}
where $Q' \triangleq \Pi \times Q$,
$AP \triangleq \{a_1,\, a_2,\, \cdots,\, a_K\}$,
$Q_0' \triangleq \Pi_0 \times Q_0$, and
$\mathcal{F}' \triangleq \Pi \times \mathcal{F}$. The transition relation
$\delta' : Q' \rightarrow 2^{Q'}$ satisfies
$\langle \pi_j,\, q_n \rangle \in \delta'(\langle \pi_i,\, q_m \rangle)$ if and
only if $(\pi_i,\, \pi_j) \in \longrightarrow_c$ and there exists
$l \in 2^{AP}$ such that $q_n \in \delta(q_m,\, l)$. The function
$W_r : Q' \times Q' \rightarrow \mathbb{R}^{+}$ is defined below.
\end{definition}

Relative to the standard product automaton in
Definition~\ref{ch02:def:product}, Definition~\ref{ch03:def:relaxpro} weakens
the automaton-side transition condition and augments the transition cost. The
weight now penalizes the distance between the observed label and the set of
labels that would enable the desired automaton transition.
To define this penalty, let
$\texttt{Eval} : 2^{AP} \rightarrow \{0,\, 1\}^{K}$ map each alphabet to its
Boolean encoding. For $l,l' \in 2^{AP}$, define
$\rho(l,\, l') \triangleq \|\boldsymbol{\nu} - \boldsymbol{\nu}'\|_{1}$,
where $\boldsymbol{\nu} = \texttt{Eval}(l)$ and
$\boldsymbol{\nu}' = \texttt{Eval}(l')$. Thus $\rho$ is the Hamming distance
between the two encodings. For any nonempty set $\chi \subseteq 2^{AP}$, define
\begin{equation}
\texttt{Dist}(l,\, \chi) \triangleq
\begin{cases}
0, & \text{if } l \in \chi, \\
\min_{l' \in \chi} \rho(l,\, l'), & \text{otherwise},
\end{cases}
\label{ch03:eq:setdist}
\end{equation}
which is well defined for $\chi \neq \emptyset$ and vanishes exactly when
$l \in \chi$. The transition weight of the relaxed product automaton is given by:
\begin{equation}
\begin{aligned}
W_r(\langle \pi_i,\, q_m \rangle,\,
\langle \pi_j,\, q_n \rangle)
\triangleq {} & w_c(\pi_i,\, \pi_j) \\
& {} + \alpha \cdot
\texttt{Dist}(L_c(\pi_i),\, \chi(q_m,\, q_n)),
\end{aligned}
\label{ch03:eq:WR}
\end{equation}
where $\langle \pi_j,\, q_n \rangle \in
\delta'(\langle \pi_i,\, q_m \rangle)$ and $\alpha \geq 0$. Here,
$\chi(q_m,\, q_n) \triangleq
\{\, l \in 2^{AP} \mid q_n \in \delta(q_m,\, l) \,\}$ collects all alphabets
that enable the transition from $q_m$ to $q_n$ in
$\mathcal{A}_{\varphi}$. By Definition~\ref{ch03:def:relaxpro}, this set is
nonempty whenever the relaxed transition exists. Thus
\eqref{ch03:eq:WR} combines implementation cost with a penalty for violating
the original automaton constraint; larger $\alpha$ gives stronger preference to
preserving the original task.

\begin{example}\label{ch03:examp:dist}
Consider the NBA introduced in Example~\ref{ch02:examp:NBA} and the transition from $q_1$
to $q_3$. In this case,
$\chi(q_1,\, q_3)
= \{\{a_2\},\, \{a_1,\, a_2\},\, \{a_2,\, a_3\},\,
\{a_1,\, a_2,\, a_3\}\}$. Therefore,
$\textup{\texttt{Dist}}(\{a_1\},\, \chi(q_1,\, q_3))
= \rho(\{a_1\},\, \{a_1,\, a_2\}) = 1$,
$\textup{\texttt{Dist}}(\{a_2\},\, \chi(q_1,\, q_3))
= \rho(\{a_2\},\, \{a_2\}) = 0$, and
$\textup{\texttt{Dist}}(\{a_2,\, a_3\},\, \chi(q_1,\, q_3))
= \rho(\{a_2,\, a_3\},\, \{a_2,\, a_3\}) = 0$.
\end{example}

\subsubsection{Balanced accepting run}\label{ch03:subsec:probStatement}

Since $\mathcal{A}_p$ has no accepting run, the search is performed on
$\mathcal{A}_r$. Different accepting runs may have different motion costs and
different specification violations, so the prefix-suffix structure from
Chapter~\ref{chap:preliminaries} is retained:
\[
R \triangleq q_0' q_1' \cdots q_{f-1}'
[\, q_f' q_{f+1}' \cdots q_n' \,]^{\omega}
= \langle \pi_0,\, q_0 \rangle \cdots
\langle \pi_{f-1},\, q_{f-1} \rangle
[\, \langle \pi_f,\, q_f \rangle
\langle \pi_{f+1},\, q_{f+1} \rangle \cdots
\langle \pi_n,\, q_n \rangle \,]^{\omega},
\]
where $q_0' = \langle \pi_0,\, q_0 \rangle \in Q_0'$ and
$q_f' = \langle \pi_f,\, q_f \rangle \in \mathcal{F}'$.
The total cost of~$R$ in the relaxed product automaton is defined as:
\begin{equation}
\texttt{Cost}(R,\, \mathcal{A}_r)
\triangleq \texttt{cost}_{\tau}
+ \alpha \cdot \texttt{dist}_{\varphi},
\label{ch03:eq:prod-cost2}
\end{equation}
where $\tau \triangleq R|_{\Pi}$ is the induced motion plan. The implementation
cost in \eqref{ch03:eq:prod-cost2} is
\begin{equation}
\texttt{cost}_{\tau} \triangleq
\Big( \sum_{i=0}^{f-1} + \gamma \sum_{i=f}^{n-1} \Big)
w_c(\pi_i,\, \pi_{i+1}),
\label{ch03:eq:cost}
\end{equation}
and the accumulated specification distance is
\begin{equation}
\texttt{dist}_{\varphi} \triangleq
\Big( \sum_{i=0}^{f-1} + \gamma \sum_{i=f}^{n-1} \Big)
\texttt{Dist}(L_c(\pi_i),\, \chi(q_i,\, q_{i+1})),
\label{ch03:eq:dist}
\end{equation}
where $\gamma \geq 0$ weights the relative importance of the transient prefix
and the recurrent suffix.

\begin{problem}\label{ch03:pro:balanced}
Find an accepting run of $\mathcal{A}_r$ that minimizes the cost in
\eqref{ch03:eq:prod-cost2}.
\end{problem}

The solution of Problem~\ref{ch03:pro:balanced} is the balanced accepting run
$R_{\textup{bal}}$ of~$\mathcal{A}_r$, with induced plan
$\tau_{\textup{bal}} \triangleq R_{\textup{bal}}|_{\Pi}$. Since NBA transitions
are represented by Boolean formulas, the distance in \eqref{ch03:eq:setdist}
can be evaluated through a binary decision diagram without enumerating
$2^{AP}$: disjunction selects the smaller branch distance, while conjunction
accumulates branch distances.


\subsubsection{Balanced plan synthesis}\label{ch03:subsec:singleFeedback}

For fixed $\alpha$ and $\gamma$, the relaxed automaton $\mathcal{A}_r$ may be
built explicitly from Definition~\ref{ch03:def:relaxpro} or generated on the
fly. Algorithm~\ref{ch02:alg:optimal} then computes
$R_{\textup{bal}}$ and the plan
$\tau_{\textup{bal}} = R_{\textup{bal}}|_{\Pi}$. Although
$\mathcal{A}_r$ admits more automaton transitions than $\mathcal{A}_p$, it does
not relax the transition relation of~$\mathcal{T}_c$; hence every balanced plan
is implementable.
Once $R_{\textup{bal}}$ is known, $\texttt{cost}_{\tau}$ and
$\texttt{dist}_{\varphi}$ follow from \eqref{ch03:eq:cost} and
\eqref{ch03:eq:dist}. A relaxed automaton~$\mathcal{A}_{\varphi}'$ can also be
obtained by adding, along $R_{\textup{bal}}$, the automaton edges needed to
support the observed labels, yielding a relaxation in the sense of
\cite{kim2012approximate}.
The parameter~$\alpha$ controls the trade-off between implementation cost and
specification violation. In practice, Algorithm~\ref{ch02:alg:optimal} can be
run for several values of $\alpha$, after which distinct balanced runs are
compared through $(\texttt{cost}_{\tau},\, \texttt{dist}_{\varphi})$ and the
selected plan is executed by Algorithm~\ref{ch02:alg:hybrid}.

\begin{lemma}\label{ch03:lem:distanceProof}
If $\textup{\texttt{dist}}_{\varphi} = 0$, then
$\tau_{\textup{bal}} \models \varphi$.
\end{lemma}

\begin{proof}
By \eqref{ch03:eq:setdist}, the function $\texttt{Dist}(\cdot,\, \cdot)$ is
nonnegative. Therefore, $\texttt{dist}_{\varphi} = 0$ implies that
$q_n \in \delta(q_m,\, L_c(\pi_i))$ for every transition
$(\langle \pi_i,\, q_m \rangle,\,
\langle \pi_j,\, q_n \rangle)$ along $R_{\textup{bal}}$. Hence every
transition of $R_{\textup{bal}}$ also satisfies the transition condition of the
standard product automaton $\mathcal{A}_p$ from
Definition~\ref{ch02:def:product}. Since $\mathcal{A}_p$ and $\mathcal{A}_r$
share the same state set, the same initial states, and the same accepting
states, $R_{\textup{bal}}$ is also an accepting run of~$\mathcal{A}_p$.
Lemma~\ref{ch02:lem:satisfy} then yields
$\tau_{\textup{bal}} \models \varphi$.
\end{proof}


\begin{theorem}\label{ch03:thm:balanced-feasible}
If $\varphi$ is feasible over $\mathcal{T}_c$, then the balanced plan
$\tau_{\textup{bal}}$ satisfies $\varphi$ whenever
$\alpha > \underline{\alpha}$, where $\underline{\alpha}$ is given by
\eqref{ch03:underalpha}.
\end{theorem}

\begin{proof}
Suppose that $\varphi$ is feasible over $\mathcal{T}_c$. Then
Algorithm~\ref{ch02:alg:optimal} returns the optimal accepting run
$R_{\textup{opt}}$ of $\mathcal{A}_p$, whose total cost is
\begin{equation}
\texttt{Cost}(R_{\textup{opt}},\, \mathcal{A}_p)
\triangleq \underline{\alpha},
\label{ch03:underalpha}
\end{equation}
under the same value of $\gamma$ as in \eqref{ch03:eq:prod-cost2}. Since every
transition of $\mathcal{A}_p$ is also a transition of $\mathcal{A}_r$,
$R_{\textup{opt}}$ is an accepting run of $\mathcal{A}_r$ as well. Moreover,
because $R_{\textup{opt}}$ satisfies the original automaton constraints
exactly, its distance term vanishes, and thus
$\texttt{Cost}(R_{\textup{opt}},\, \mathcal{A}_r)
= \texttt{Cost}(R_{\textup{opt}},\, \mathcal{A}_p)
= \underline{\alpha}$.
Assume, for contradiction, that $\tau_{\textup{bal}} \not\models \varphi$.
Then \eqref{ch03:eq:dist} implies $\texttt{dist}_{\varphi} \geq 1$. Therefore,
by \eqref{ch03:eq:prod-cost2},
$\texttt{Cost}(R_{\textup{bal}},\, \mathcal{A}_r)
= \texttt{cost}_{\tau} + \alpha \cdot \texttt{dist}_{\varphi}
\geq \alpha$. If $\alpha > \underline{\alpha}$, then
$\texttt{Cost}(R_{\textup{bal}},\, \mathcal{A}_r)
> \underline{\alpha}
= \texttt{Cost}(R_{\textup{opt}},\, \mathcal{A}_r)$, which contradicts the
optimality of $R_{\textup{bal}}$ for Problem~\ref{ch03:pro:balanced}. Hence
$\tau_{\textup{bal}} \models \varphi$ holds.
\end{proof}



\begin{example}\label{ch03:examp:singleRS}
Consider a robot that starts from region~$\pi_0$ and should eventually reach~$\pi_3$ and remain there, while
avoiding every region that satisfies $a_2$ or $a_3$. By varying $\alpha$ with
$\gamma = 5$, three distinct balanced plans are obtained:
(I) When the penalty on violating~$\varphi$ is small,
$\mathcal{A}_{\varphi}$ is revised by adding $q_1$ to
$\delta(q_0,\, \emptyset)$ and $q_1$ to $\delta(q_1,\, \emptyset)$. The
balanced plan is then $[\pi_0]^{\omega}$, indicated by the black hexagram, with
$\textup{\texttt{cost}}_{\tau} = 30$ and
$\textup{\texttt{dist}}_{\varphi} = 6$.
(II) When the penalty is increased, $\mathcal{A}_{\varphi}$ is revised by adding
$q_1$ to $\delta(q_0,\, \{a_2,\, a_3\})$. The resulting balanced plan is
$\pi_0\, \pi_1\, [\pi_3]^{\omega}$, indicated by the blue square, with
$\textup{\texttt{cost}}_{\tau} = 65$ and
$\textup{\texttt{dist}}_{\varphi} = 2$.
(III) When the penalty becomes sufficiently large,
$\mathcal{A}_{\varphi}$ is revised by adding $q_1$ to
$\delta(q_0,\, \{a_2\})$. The balanced plan becomes
$\pi_0\, \pi_2\, [\pi_3]^{\omega}$, indicated by the cyan triangle, with
$\textup{\texttt{cost}}_{\tau} = 85$ and
$\textup{\texttt{dist}}_{\varphi} = 1$. This case shows that the method
prefers to pass through~$\pi_2$, which satisfies only~$a_2$, rather than
through $\pi_1$, which satisfies both~$a_2$ and~$a_3$.
\end{example}












%%%%%%%%%%%%%%%%%%%%%%%%%%%%%%%%%%%%%%%%%%%%%%%%%%%
\subsection{Soft and hard specifications}\label{ch03:sec:soft-hard}

The previous subsection allowed the entire specification to be relaxed. Many
applications instead separate safety-critical requirements, which must never be
violated, from performance or achievement objectives, which may be relaxed under
partial workspace knowledge. This subsection develops a synthesis framework for
that hard--soft structure.

\subsubsection{Problem formulation}\label{ch03:subsec:formulation}

The robot transition system is still denoted by $\mathcal{T}_c$, as in
Definition~\ref{ch02:def:FTS}. The task specification remains an LTL formula
over $AP$, but now has the form
\begin{equation}
\varphi \triangleq \varphi^{\textup{soft}} \wedge \varphi^{\textup{hard}},
\label{ch03:eq:spec}
\end{equation}
where $\varphi^{\textup{hard}}$ and $\varphi^{\textup{soft}}$ are the hard and
soft components. Hard requirements include constraints such as obstacle
avoidance or recurrent charging; soft requirements encode objectives whose
feasibility may depend on incomplete workspace information. Thus relaxation is
allowed only for the soft component.

\begin{problem}
Given the specification in \eqref{ch03:eq:spec}, synthesize a motion and task
plan such that $\varphi^{\textup{hard}}$ is satisfied fully, while
$\varphi^{\textup{soft}}$ is satisfied as much as possible.
\end{problem}

\subsubsection{Safety-ensured product automaton}
\label{ch03:subsec:intersection}

The decomposition in \eqref{ch03:eq:spec} leads naturally to three notions of
correctness.

\begin{definition}\label{ch03:def:vs}
A finite or infinite path $\tau$ of $\mathcal{T}_c$ is called \emph{valid}
if every consecutive pair of states along $\tau$ belongs to
$\longrightarrow_c$. An infinite valid path $\tau$ is called \emph{safe} if
$\tau \models \varphi^{\textup{hard}}$, and \emph{satisfying} if
$\tau \models \varphi$.
\end{definition}

Let
$\mathcal{A}^{\textup{hard}}
= (Q_1,\, 2^{AP},\, \delta_1,\, Q_{1,0},\, \mathcal{F}_1)$ and
$\mathcal{A}^{\textup{soft}}
= (Q_2,\, 2^{AP},\, \delta_2,\, Q_{2,0},\, \mathcal{F}_2)$ be the NBAs for
$\varphi^{\textup{hard}}$ and $\varphi^{\textup{soft}}$, respectively. Define
$\chi_1(q_1,\, \check{q}_1) \triangleq
\{\, l \in 2^{AP} \mid \check{q}_1 \in \delta_1(q_1,\, l) \,\}$ and
$\chi_2(q_2,\, \check{q}_2) \triangleq
\{\, l \in 2^{AP} \mid \check{q}_2 \in \delta_2(q_2,\, l) \,\}$ for the hard
and soft automata. The hard component is enforced exactly, while the soft
component is relaxed in the following intersection automaton.

\begin{definition}\label{ch03:def:conj}
The relaxed intersection of $\mathcal{A}^{\textup{hard}}$ and
$\mathcal{A}^{\textup{soft}}$ is
\begin{equation}
\tilde{\mathcal{A}}_{\varphi}
\triangleq
(Q,\, 2^{AP},\, \delta,\, Q_0,\, \mathcal{F}),
\label{ch03:eq:intersect}
\end{equation}
where
$Q \triangleq Q_1 \times Q_2 \times \{1,\, 2\}$,
$Q_0 \triangleq Q_{1,0} \times Q_{2,0} \times \{1\}$, and
$\mathcal{F} \triangleq \mathcal{F}_1 \times Q_2 \times \{1\}$. The transition
relation $\delta : Q \times 2^{AP} \rightarrow 2^Q$ satisfies
$\langle \check{q}_1,\, \check{q}_2,\, \check{t} \rangle \in
\delta(\langle q_1,\, q_2,\, t \rangle,\, l)$ if and only if
$l \in \chi_1(q_1,\, \check{q}_1)$,
$\chi_2(q_2,\, \check{q}_2) \neq \emptyset$, and either
$q_t \notin \mathcal{F}_t$ with $\check{t} = t$, or
$q_t \in \mathcal{F}_t$ with
$\check{t} = \textup{\texttt{mod}}(t,\, 2) + 1$.
\end{definition}

Definition~\ref{ch03:def:conj} relaxes only the soft component: the actual
alphabet must satisfy the hard transition, whereas the soft transition must
only be enabled by some alphabet. The auxiliary variable $t \in \{1,\, 2\}$
alternates the acceptance obligation as in the standard NBA intersection
\cite{baier2008principles}. Denote by $r|_{Q_1}$ and $r|_{Q_2}$ the projections
of a run~$r$ onto the hard and soft automata.

\begin{algorithm}[t]
\DontPrintSemicolon
\caption{Construction of the relaxed intersection, \texttt{RelaxInt}()}
\label{ch03:alg:relaxedint}
\KwIn{$\mathcal{A}^{\textup{soft}}$, $\mathcal{A}^{\textup{hard}}$}
\KwOut{$\tilde{\mathcal{A}}_{\varphi}$}

\ForEach{$q_1 \in Q_1$, $q_2 \in Q_2$, $t \in \{1,\,2\}$}{
  $q_m = \langle q_1,\, q_2,\, t \rangle \in Q$\;

  \If{$q_1 \in Q_{1,0}$, $q_2 \in Q_{2,0}$, and $t = 1$}{
    $q_m \in Q_0$\;
  }

  \If{$q_1 \in \mathcal{F}_1$ and $t = 1$}{
    $q_m \in \mathcal{F}$\;
  }

  \ForEach{$\check{q}_1 \in \textup{\texttt{Post}}(q_1)$,
           $\check{q}_2 \in \textup{\texttt{Post}}(q_2)$,
           $\check{t} \in \{1,\,2\}$}{
    $q_n = \langle \check{q}_1,\, \check{q}_2,\, \check{t} \rangle \in Q$\;

    \If{$(q_1 \notin \mathcal{F}_1$ and $t = 1$ and $\check{t} = 1)$ or
        $(q_2 \notin \mathcal{F}_2$ and $t = 2$ and $\check{t} = 2)$ or
        $(q_1 \in \mathcal{F}_1$ and $t = 1$ and $\check{t} = 2)$ or
        $(q_2 \in \mathcal{F}_2$ and $t = 2$ and $\check{t} = 1)$}{
      $q_n \in \delta(q_m,\, l)$, for all $l \in \chi_1(q_1,\, \check{q}_1)$\;
    }
  }
}
\textbf{return} $\tilde{\mathcal{A}}_{\varphi}$
\end{algorithm}

\begin{theorem}\label{ch03:theo:saferun}
Given an accepting run $r$ of $\tilde{\mathcal{A}}_{\varphi}$, the projection
$r|_{Q_1}$ is an accepting run of $\mathcal{A}^{\textup{hard}}$. Moreover,
$\mathcal{L}_{\omega}(\tilde{\mathcal{A}}_{\varphi}) \subseteq
\mathcal{L}_{\omega}(\mathcal{A}^{\textup{hard}})$.
\end{theorem}

\begin{proof}
Since $r$ is accepting, it visits the set
$\mathcal{F} = \mathcal{F}_1 \times Q_2 \times \{1\}$ infinitely often.
Therefore, the projection $r|_{Q_1}$ visits $\mathcal{F}_1$ infinitely often.
In addition, every transition of $r$ satisfies
$l \in \chi_1(q_1,\, \check{q}_1)$ by
Definition~\ref{ch03:def:conj}. Hence every projected transition of $r|_{Q_1}$
is a valid transition of $\mathcal{A}^{\textup{hard}}$. It follows that
$r|_{Q_1}$ is an accepting run of $\mathcal{A}^{\textup{hard}}$.
For the language inclusion, let
$\sigma \in \mathcal{L}_{\omega}(\tilde{\mathcal{A}}_{\varphi})$. Then
$\sigma$ induces an accepting run $r_{\sigma}$ of
$\tilde{\mathcal{A}}_{\varphi}$. By the first part of the proof,
$r_{\sigma}|_{Q_1}$ is an accepting run of
$\mathcal{A}^{\textup{hard}}$. Therefore,
$\sigma \in \mathcal{L}_{\omega}(\mathcal{A}^{\textup{hard}})$, which proves
the desired inclusion.
\end{proof}

To preserve the hard specification while still optimizing the soft one, combine
$\tilde{\mathcal{A}}_{\varphi}$ with the workspace transition system in a
relaxed product automaton.

\begin{definition}\label{ch03:def:safe}
The safety-ensured relaxed product B\"uchi automaton is defined as:
\begin{equation}
\tilde{\mathcal{A}}_{r}
\triangleq
\mathcal{T}_c \times \tilde{\mathcal{A}}_{\varphi}
=
(Q',\, 2^{AP},\, \delta',\, Q_0',\, \mathcal{F}',\, W_r),
\label{ch03:eq:safe-product}
\end{equation}
where $Q' \triangleq \Pi \times Q$,
$Q_0' \triangleq \Pi_0 \times Q_0$, and
$\mathcal{F}' \triangleq \Pi \times \mathcal{F}$. The transition relation
$\delta' : Q' \rightarrow 2^{Q'}$ satisfies
$\langle \pi_j,\, q_n \rangle \in \delta'(\langle \pi_i,\, q_m \rangle)$ if and
only if $(\pi_i,\, \pi_j) \in \longrightarrow_c$ and
$q_n \in \delta(q_m,\, L_c(\pi_i))$. The weight function
$W_r : Q' \times Q' \rightarrow \mathbb{R}^{+}$ is given by:
\begin{equation}
\begin{aligned}
W_r(\langle \pi_i,\, q_m \rangle,\,
\langle \pi_j,\, q_n \rangle)
\triangleq {} & w_c(\pi_i,\, \pi_j) \\
& {} + \alpha \cdot
\textup{\texttt{Dist}}(L_c(\pi_i),\, \chi_2(q_2,\, \check{q}_2)),
\end{aligned}
\label{ch03:eq:weight-soft}
\end{equation}
where $q_m = \langle q_1,\, q_2,\, t \rangle$,
$q_n = \langle \check{q}_1,\, \check{q}_2,\, \check{t} \rangle$, and
$\alpha \geq 0$.
\end{definition}

The weight in \eqref{ch03:eq:weight-soft} combines implementation cost with
the violation distance of the soft specification. The hard component is absent
from the weight because it is enforced directly by the transition relation of
$\tilde{\mathcal{A}}_{\varphi}$.

\begin{theorem}\label{ch03:theo:safety}
Assume that $R$ is an accepting run of $\tilde{\mathcal{A}}_{r}$. Then its
projection $\tau \triangleq R|_{\Pi}$ is both valid and safe in the sense of
Definition~\ref{ch03:def:vs}.
\end{theorem}

\begin{proof}
The path $\tau$ is valid because every transition in $\delta'$ projects onto a
transition in $\longrightarrow_c$. Since $R$ is an accepting run of
$\tilde{\mathcal{A}}_{r}
= \mathcal{T}_c \times \tilde{\mathcal{A}}_{\varphi}$, the trace of $\tau$ is
accepted by $\tilde{\mathcal{A}}_{\varphi}$. By
Theorem~\ref{ch03:theo:saferun},
$\texttt{trace}(\tau) \in
\mathcal{L}_{\omega}(\mathcal{A}^{\textup{hard}})$. Then, by
Lemma~\ref{ch02:lem:LTL-BA},
$\texttt{trace}(\tau) \in \texttt{Words}(\varphi^{\textup{hard}})$, which is
equivalent to $\tau \models \varphi^{\textup{hard}}$. Hence, the resulting plan $\tau$ is safe.
\end{proof}

As in Section~\ref{ch03:sec:infeasible}, accepting runs are evaluated in the
prefix-suffix form of \eqref{ch02:eq:prefix-suffix}. Their total cost is
defined by
\begin{equation}
\begin{aligned}
\texttt{Cost}(R,\, \tilde{\mathcal{A}}_{r})
&\triangleq
\texttt{cost}_{\tau}
+ \alpha \cdot \texttt{dist}_{\varphi^{\textup{soft}}}, \\
\texttt{dist}_{\varphi^{\textup{soft}}}
&\triangleq
\Big( \sum_{i=0}^{f-1} + \gamma \sum_{i=f}^{n-1} \Big)
\texttt{Dist}(L_c(\pi_i),\,
\chi_2(q_i'|_{Q_2},\, q_{i+1}'|_{Q_2})),
\end{aligned}
\label{ch03:eq:prod-cost3}
\end{equation}
where $\tau \triangleq R|_{\Pi}$, $\texttt{cost}_{\tau}$ is given by
\eqref{ch03:eq:cost}, and $q_i'|_{Q_2}$ denotes the projection of the product
state $q_i'$ onto the soft-automaton state set $Q_2$.

\begin{problem}\label{ch03:pro:safe}
Find an accepting run of $\tilde{\mathcal{A}}_{r}$ that minimizes the cost in
\eqref{ch03:eq:prod-cost3}.
\end{problem}

\subsubsection{Safe plan synthesis}\label{ch03:subsec:safe-plan-synthesis}

The solution of Problem~\ref{ch03:pro:safe} is the safe accepting run
$R_{\textup{safe}}$ of $\tilde{\mathcal{A}}_{r}$, with induced plan
$\tau_{\textup{safe}} \triangleq R_{\textup{safe}}|_{\Pi}$. For fixed
$\alpha$ and $\gamma$, $\tilde{\mathcal{A}}_{r}$ can be built explicitly or on
the fly. Algorithm~\ref{ch02:alg:optimal} then computes
$R_{\textup{safe}}$, and Theorem~\ref{ch03:theo:safety} guarantees that
$\tau_{\textup{safe}}$ is valid and safe. As before, varying $\alpha$ yields
candidate runs that can be compared through
$(\texttt{cost}_{\tau},\, \texttt{dist}_{\varphi^{\textup{soft}}})$ before
implementation by Algorithm~\ref{ch02:alg:hybrid}.

\begin{lemma}\label{ch03:lem:distanceProofsafe}
If $\textup{\texttt{dist}}_{\varphi^{\textup{soft}}} = 0$, then
$\tau_{\textup{safe}} \models \varphi$.
\end{lemma}

\begin{proof}
Let $\mathcal{A}_{\varphi}$ denote the standard, unrelaxed intersection of
$\mathcal{A}^{\textup{hard}}$ and $\mathcal{A}^{\textup{soft}}$. Since
$\texttt{dist}_{\varphi^{\textup{soft}}} = 0$, every transition along
$R_{\textup{safe}}$ satisfies the soft component exactly. The hard component is
already enforced by the transition relation of
$\tilde{\mathcal{A}}_{\varphi}$. Therefore, $R_{\textup{safe}}$ is also an
accepting run of the standard product automaton
$\mathcal{T}_c \times \mathcal{A}_{\varphi}$. By
Lemma~\ref{ch02:lem:satisfy}, the plan
$\tau_{\textup{safe}}$ satisfies $\varphi$.
\end{proof}














%%%%%%%%%%%%%%%%%%%%%%%%%%%%%%%%%%%%%%%%%%%%%%%%%%%%
%%%%%%%%%%%%%%%%%%%%%%%%%%%%%%%%%%%%%%%%%%%%%%%%%%%
%%%%%%%%%%%%%%%%%%%%%%%%%%%%%%%%%%%%%%%%%%%%%%%%%%%%
\subsection{Online sensing and plan adaptation}
\label{ch03:sec:partial}

In practice, the transition-system model may become inconsistent with the
actual workspace during execution. A plan synthesized off-line can therefore
become invalid, unsafe, or suboptimal after new information is acquired. This
motivates an online framework in which the robot monitors execution, updates the
workspace model, and revises the current plan while preserving validity and
safety. Let
\begin{equation}\label{ch03:eq:t-FTS}
\mathcal{T}_c^t\triangleq (\Pi,\, \longrightarrow_c^t,\, \Pi_0,\,AP,\, L_c^t,\, w_c^t),
\end{equation}
denote the transition system available at time $t \geq 0$. The region set
$\Pi$, the initial-state set $\Pi_0$, and the proposition set $AP$ are taken
to be fixed. Only the transition relation, the labeling function, and the
transition weights are updated over time. The task specification is also fixed
during execution and retains the decomposition
$\varphi \triangleq
\varphi^{\textup{soft}} \wedge \varphi^{\textup{hard}}$
introduced in \eqref{ch03:eq:spec}. Let
$\tilde{\mathcal{A}}_{\varphi}$ denote the relaxed intersection automaton from
Definition~\ref{ch03:def:conj}.

\begin{problem}
Assume that the workspace is only partially known. The following questions
arise: (I) how should the robot's sensing information be modeled;
(II) how should the transition system be updated accordingly;
and (III) how can the motion and task plan be kept valid and safe throughout execution?
\end{problem}

\subsubsection{Initial plan synthesis}

At $t=0$, the initial plan is synthesized by
Section~\ref{ch03:sec:soft-hard} from $\mathcal{T}_c^0$ and~$\varphi$. For
$t \geq 0$, let $R^t$ be the current safe accepting run,
$\tau^t \triangleq R^t|_{\Pi}$ the corresponding plan, and
$\tilde{\mathcal{A}}_r^t$ the associated product automaton, generated on the
fly. Even if the soft component is initially infeasible,
Theorem~\ref{ch03:theo:safety} guarantees that $\tau^0$ is valid and safe with
respect to $\mathcal{T}_c^0$.

\subsubsection{Real-time knowledge update}
\label{ch03:subsec:knowledge-update}

The robot acquires workspace information through sensing during execution.
Only locally observed information is used here; inter-agent information is
introduced later. Let $\textbf{Sense}^t$ denote the sensing information obtained
at time $t \geq 0$, either upon reaching a region or while traversing a
transition:
\begin{equation}
\textbf{Sense}^t \triangleq
\{\, ((\pi,\, S,\, S_{\neg}),\, E,\, E_{\neg}) \,\},
\label{ch03:eq:sense-info}
\end{equation}
where $\pi \in \Pi$ is an observed region, $S$ and $S_{\neg}$ contain
propositions confirmed true and false at~$\pi$, respectively,
$(\pi_i,\, \pi_j,\, w) \in E$ adds or reweights a transition, and
$(\pi_i,\, \pi_j) \in E_{\neg}$ removes one. Thus $\textbf{Sense}^t$ encodes
the symbolic workspace information revealed at time~$t$. The extraction of
\eqref{ch03:eq:sense-info} from raw sensor data depends on the sensing modality
and is not needed here.

\subsubsection{Transition system update}\label{ch03:subsubsec:updateT}

At time~$t$, the robot may receive new information through
$\textbf{Sense}^t$. This information is used to update the current
transition-system model. Let $\mathcal{T}_c^{t^-}$ and
$\mathcal{T}_c^{t^+}$ denote the transition systems immediately before and
after the update at time~$t$, respectively.
If $(\pi,\, S,\, S_{\neg})$ is obtained from \eqref{ch03:eq:sense-info}, then
the labeling function is updated according to
$L_c^{t^+}(\pi) \triangleq
\big(L_c^{t^-}(\pi) \cup S\big) \setminus S_{\neg}$.
Similarly, each element of~$E$ adds a transition or updates its weight, whereas
each element of~$E_{\neg}$ removes a transition. For convenience,
$\widetilde{\Pi} \subseteq \Pi$ denotes the set of regions whose labeling
function changes during the update, and
\begin{equation}
\widehat{\Pi} \triangleq
\widetilde{\Pi} \cup
\{\, \pi_i \mid (\pi_i,\, \pi_j,\, w) \in E
\text{ or } (\pi_i,\, \pi_j) \in E_{\neg} \,\},
\label{ch03:eq:widehat}
\end{equation}
denotes the set of regions whose outgoing adjacency information must be
reconstructed. If $\textbf{Sense}^t$ is empty, then
$\mathcal{T}_c^{t^+} = \mathcal{T}_c^{t^-}$.

\subsubsection{Safety-ensured plan revision}\label{ch03:subsec:realtime}

Because $\mathcal{T}_c^t$ may change during execution, the current plan must
be re-evaluated for validity, safety, and optimality. The update first modifies
the product automaton representation and then decides whether the current run
can be retained, locally repaired, or replaced.

(I) \emph{Product automaton update}.
Let $\tilde{\mathcal{A}}_r^{t^-}$ and $\tilde{\mathcal{A}}_r^{t^+}$ denote the
product automata before and after the update. Reconstructing
$\tilde{\mathcal{A}}_r^{t^+}$ globally, or re-evaluating every affected
transition as in~\cite{guo2013revising}, is often unnecessary. Instead, update
information is incorporated lazily into the adjacency routine: product states
whose workspace component lies in $\widehat{\Pi}$ have their outgoing
transitions rebuilt only when revisited. The remaining questions are whether
$\tau^t$ is still valid and safe, and how it should be revised if not.

(II) \emph{Verifying validity and safety}.
The prefix-suffix representation of $R^t$ contains finitely many distinct
transitions. Let $\texttt{Edge}(R^t)$ be the set of transitions appearing in
that representation. Scanning this set determines whether the current plan
remains valid and safe.

\begin{definition}\label{ch03:def:vsTS}
Given the updated transition system $\mathcal{T}_c^{t^+}$, a transition
$(\langle \pi_i,\, q_m \rangle,\,
\langle \pi_j,\, q_n \rangle) \in \textup{\texttt{Edge}}(R^t)$ is called
\emph{invalid} if $(\pi_i,\, \pi_j) \notin \longrightarrow_c^{t^+}$, and
\emph{unsafe} if
$L_c^{t^+}(\pi_i) \notin \chi_1(q_m|_{Q_1},\, q_n|_{Q_1})$.
\end{definition}

Let $\Xi$ and $\aleph$ denote the sets of invalid and unsafe transitions of
$R^t$, respectively. These sets are obtained by scanning the transitions in
$\texttt{Edge}(R^t)$ and checking whether a transition has been removed from
$\mathcal{T}_c^{t^+}$ or whether an updated label violates the hard-automaton
constraint.

\begin{theorem}\label{ch03:theo:validity}
Assume that $R^t$ is an accepting run of $\tilde{\mathcal{A}}_r^{t^-}$ and that
$\mathcal{T}_c^{t^-}$ is updated to $\mathcal{T}_c^{t^+}$. Let $\Xi$ and
$\aleph$ be the invalid and unsafe transitions above. Then the
following statements hold:
(I) $\tau^t$ remains valid if and only if $\Xi = \emptyset$;
(II) $\tau^t$ remains safe if $\Xi = \emptyset$ and $\aleph = \emptyset$.
\end{theorem}

\begin{proof}
Since $R^t$ is an accepting run of $\tilde{\mathcal{A}}_r^{t^-}$,
Theorem~\ref{ch03:theo:safety} implies that $\tau^t$ is valid and safe for
$\mathcal{T}_c^{t^-}$. Statement~(i) follows directly from the definition of
$\Xi$: the plan remains valid precisely when none of its transitions has been
removed from the updated transition relation.
For statement~(ii), assume that $\Xi = \emptyset$ and $\aleph = \emptyset$.
Then every transition of $R^t$ remains valid in $\mathcal{T}_c^{t^+}$ and
still satisfies the hard-automaton constraint under the updated labeling
function. Hence $R^t$ remains an accepting run of
$\tilde{\mathcal{A}}_r^{t^+}$, and Theorem~\ref{ch03:theo:safety} yields that
$\tau^t$ remains safe.
\end{proof}

(III) \emph{Plan revision}:
If either $\Xi$ or $\aleph$ is nonempty, the current plan may no longer be
executable or safe. Replanning from the current region alone can yield a valid
future plan, but it may be incompatible with the product-state history induced
by the executed trajectory. Let $\tau_{\textup{past}}$ and
$R_{\textup{past}}$ denote the executed trajectory and its recorded product run,
and let
$\tau_{\textup{comp}} \triangleq
\tau_{\textup{past}} + \tau_{\textup{new}}$. Since the B\"uchi automaton is
nondeterministic, several product states may correspond to the same current
workspace region, which leads to the following compatibility problem.

\begin{problem}\label{ch03:prob:concatenate}
Assume that $\Xi$ or $\aleph$ is nonempty. Given the executed past trajectory
$\tau_{\textup{past}}$, find a new plan $\tau_{\textup{new}}$ such that the
complete trajectory
$\tau_{\textup{comp}} = \tau_{\textup{past}} + \tau_{\textup{new}}$
is valid and safe.
\end{problem}

To formalize the compatibility constraint, let
$\bar{\tau}_{\textup{past}} \triangleq
\tau_{\textup{past}}[1:(|\tau_{\textup{past}}|-1)]$, that is, the executed
trajectory without its last state. Because $\tilde{\mathcal{A}}_{\varphi}$ is
nondeterministic, the trace of $\bar{\tau}_{\textup{past}}$ may induce several
finite runs in $\tilde{\mathcal{A}}_{\varphi}$. Denote this set by
$r_{\tau_{\textup{past}}} \triangleq
\{\, r_{\bar{\sigma}}
\mid \bar{\sigma} = \texttt{trace}(\bar{\tau}_{\textup{past}}) \,\}$.
The corresponding set of finite runs in the updated product
automaton is
$R_{\tau_{\textup{past}}} \triangleq
\{\, R \mid R|_{\Pi} = \tau_{\textup{past}},\;
R|_{Q} \in r_{\tau_{\textup{past}}} \,\}$.
Accordingly, the product states consistent with the executed past is
\begin{equation}
Q'_{\tau_{\textup{past}}} \triangleq
\{\, \texttt{last}(R) \mid R \in R_{\tau_{\textup{past}}} \,\},
\label{ch03:eq:set-past}
\end{equation}
where $\texttt{last}(R)$ is the last state of the finite run~$R$. By
construction, $R_{\textup{past}} \in R_{\tau_{\textup{past}}}$ and the recorded
current product state belongs to $Q'_{\tau_{\textup{past}}}$. The revision
algorithm first attempts to repair the current run by inserting a bridging
segment around each invalid or unsafe transition, following the spirit of
\cite{guo2013revising}. If no bridge exists, it calls
\texttt{OptRun}$(\cdot)$ with initial-state set
$Q'_{\tau_{\textup{past}}}$ rather than $Q_0'$, ensuring compatibility with the
executed past.

\begin{algorithm}[t]
\caption{Revise the current plan, \texttt{Revise}()}
\label{ch03:alg:revise}
\DontPrintSemicolon
\KwIn{$\Xi$, $\aleph$, $R^{t^-}$, $\tilde{\mathcal{A}}_{r}^{t^+}$,
      $Q'_{\tau_{\textup{past}}}$}
\KwOut{$R^{t^+}$}

\ForAll{$(q'_s,\, q'_{s+1}) \in (\Xi \cup \aleph)$}{
  \If{$(q'_s,\, q'_{s+1}) \in R^{t^-}_{\textup{pre}}$}{
    $bridge = \texttt{DijksTarget}(
    \tilde{\mathcal{A}}_{r}^{t^+},\, q'_{s-1},\,
    \texttt{tail}(q'_s,\, R^{t^-}_{\textup{pre}}))$\;

    \If{$bridge \neq \emptyset$}{
      $R^{t^-}_{\textup{pre}} =
      \texttt{head}(q'_{s-1},\, R^{t^-}_{\textup{pre}})
      + bridge
      + \texttt{tail}(\texttt{last}(bridge),\, R^{t^-}_{\textup{pre}})$\;
    }
    \Else{
      $R^{t^+} =
      \texttt{OptRun}(\tilde{\mathcal{A}}_{r}^{t^+},
      Q'_{\tau_{\textup{past}}})$\;
      \textbf{return} $R^{t^+}$\;
    }
  }

  \If{$(q'_s,\, q'_{s+1}) \in R^{t^-}_{\textup{suf}}$}{
    Apply the same procedure to $R^{t^-}_{\textup{suf}}$\;
  }
}
\textbf{return}
$R^{t^+} = \langle R^{t^-}_{\textup{pre}},\, R^{t^-}_{\textup{suf}} \rangle$\;
\end{algorithm}

\begin{theorem}\label{ch03:theo:new-safe}
Assume that there exists an accepting run of
$\tilde{\mathcal{A}}_{r}^{t^+}$ whose projection onto $\Pi$ extends the
executed past trajectory $\tau_{\textup{past}}$. Then
Algorithm~\ref{ch03:alg:revise} returns a continuation
$\tau_{\textup{new}}$ such that
$\tau_{\textup{comp}} = \tau_{\textup{past}} + \tau_{\textup{new}}$
is valid and safe.
\end{theorem}

\begin{proof}
If a compatible accepting continuation exists, then there is an accepting run
of $\tilde{\mathcal{A}}_{r}^{t^+}$ whose initial state belongs to
$Q'_{\tau_{\textup{past}}}$. Algorithm~\ref{ch03:alg:revise} first attempts to
repair the current run locally by inserting bridging segments around the
transitions in $\Xi \cup \aleph$. Whenever such a bridge is found, the modified
run remains consistent with the unchanged portion of the original accepting
run.
If no bridge can be found for some transition, the algorithm calls
\texttt{OptRun}$(\cdot)$ with initial-state set
$Q'_{\tau_{\textup{past}}}$. By assumption, a compatible accepting run exists,
so this search succeeds and returns one such run. In either case, the resulting
run is an accepting run of $\tilde{\mathcal{A}}_{r}^{t^+}$ whose projection
extends $\tau_{\textup{past}}$. Theorem~\ref{ch03:theo:safety} then implies
that the complete trajectory $\tau_{\textup{comp}}$ is valid and safe.
\end{proof}

Algorithm~\ref{ch03:alg:scheme} later combines the revision procedure of
Algorithm~\ref{ch03:alg:revise} with the optimal search routine of
Algorithm~\ref{ch02:alg:optimal}; additional implementation details appear
in~\cite{guo2014cooperative}. The hybrid controller then receives the updated
next target region.

%%%%%%%%%%%%%%%%%%%%%%%%%%%%%%%%%%%%%%%%%%%%%
%%%%%%%%%%%%%%%%%%%%%%%%%%%%%%%%%%%%%%%%%%%%%











\subsection{Collaborative knowledge transfer}
\label{ch03:sec:multi-partial}

The previous subsection treated online sensing and adaptation for one robot.
In a multi-robot setting, each robot may keep its own local task and plan, while
workspace information sensed by one robot can improve the models used by
others. When tasks are specification-level independent, this exchange does not
require a centralized team product automaton. This subsection develops the
corresponding knowledge-transfer mechanism.
Consider a team of autonomous robots indexed by
$i \in \mathcal{N} \triangleq \{1,\, 2,\, \cdots,\, N\}$, all operating in the
same partially known workspace~$\mathcal{W}$. The motion of robot~$i$ is
abstracted by the weighted transition system defined as:
\begin{equation}
\mathcal{T}_i^t \triangleq
(\Pi_i,\, \longrightarrow_i^t,\, \Pi_{i,0},\, AP_i,\, L_i^t,\, W_i^t),
\label{ch03:eq:Tk}
\end{equation}
which is defined in the same way as \eqref{ch03:eq:t-FTS}. The superscript
$t \geq 0$ indicates that the model may be updated during execution. Each
robot~$i$ is assigned a local task of the form
$\varphi_i \triangleq
\varphi_i^{\textup{soft}} \wedge \varphi_i^{\textup{hard}}$,
where the hard and soft components have the same interpretation as in
\eqref{ch03:eq:spec}. The local tasks are assumed to be independent in the
sense that $\varphi_i$ depends only on propositions in~$AP_i$.

\subsubsection{Initial synthesis}

At time $t = 0$, each robot~$i \in \mathcal{N}$ synthesizes its own initial
plan by the framework of Section~\ref{ch03:sec:soft-hard}, using
$\mathcal{T}_i^0$ and~$\varphi_i$. Let
$\tilde{\mathcal{A}}_{\varphi_i}$ denote the relaxed intersection automaton
associated with~$\varphi_i$, and let
$\tilde{\mathcal{A}}_{r,i}^t$ denote the corresponding safety-ensured relaxed
product automaton at time~$t$. The resulting safe accepting run and induced
safe plan are denoted by $R_i^t$ and $\tau_i^t \triangleq R_i^t|_{\Pi_i}$,
respectively. Even when the soft component of $\varphi_i$ is initially
infeasible, Theorem~\ref{ch03:theo:safety} guarantees that $\tau_i^0$ is valid
and safe with respect to $\mathcal{T}_i^0$.

\subsubsection{Knowledge transfer protocol}
\label{ch03:subsec:know-transfer}

Knowledge transfer uses the team communication network. Under fixed-topology
communication, robot~$i$ has neighbor set $\mathcal{N}_i$ and
$\mathcal{N}_i^t \triangleq \mathcal{N}_i$. Under range-limited communication,
robot~$i$ can send messages to robot~$j$ only when
$\|x_j(t) - x_i(t)\| \leq C_i$, so
$\mathcal{N}_i^t \triangleq
\{\, j \in \mathcal{N} \mid \|x_j(t) - x_i(t)\| \leq C_i \,\}$. The notation
$\mathcal{N}_i^t$ covers both cases.
Robot~$i$ is interested only in propositions that appear in its own local task,
namely the set $\varphi_i|_{AP_i}$. To reduce communication load, a
subscriber-publisher mechanism is adopted. When robot~$j$ communicates with
robot~$i \in \mathcal{N}_j^t$ for the first time, robot~$j$ subscribes to
task-relevant information held by robot~$i$ by sending the request:
\begin{equation}
\textbf{Request}_{j,i}^t \triangleq
(j,\, \varphi_j|_{AP_j},\, i),
\label{ch03:eq:robot-request}
\end{equation}
which records the subscriber, the relevant propositions, and the intended
publisher. Each robot stores received subscriptions and sends each request at
most once. When robot~$i$ obtains
$((\pi,\, S,\, S_{\neg}),\, E,\, E_{\neg}) \in \textbf{Sense}_i^t$, it checks
whether $S \cup S_{\neg}$ contains propositions relevant to any reachable
subscriber. If so, it sends
\begin{equation}
\textbf{Reply}_{i,j}^t \triangleq
\{\, (\pi,\, S',\, S_{\neg}') \,\},
\label{ch03:eq:robot-reply}
\end{equation}
where
$S' \triangleq S \cap (\varphi_j|_{AP_j})$ and
$S_{\neg}' \triangleq S_{\neg} \cap (\varphi_j|_{AP_j})$.
Only task-relevant propositions are included in the reply, so every
message carries information that may affect robot~$j$'s local plan.

\begin{algorithm}[t]
\caption{Transfer knowledge to subscribed robots, \texttt{TranKnow}()}
\label{ch03:alg:transferAlgo}
\DontPrintSemicolon
\KwIn{$\textbf{Sense}_i^t$, subscriber list $\mathcal{S}_i^t$}
\KwOut{$\textbf{Reply}_{i,j}^t$ for reachable subscribers $j$}

Initialize $\textbf{Reply}_{i,j}^t = \emptyset$ for all $j \in \mathcal{N}_i^t$\;

\ForAll{$(j,\, \varphi_j|_{AP_j},\, i) \in \mathcal{S}_i^t$}{
  \If{$j \in \mathcal{N}_i^t$}{
    \ForAll{$((\pi,\, S,\, S_{\neg}),\, E,\, E_{\neg})
            \in \textbf{Sense}_i^t$}{
      $S' = S \cap (\varphi_j|_{AP_j})$\;
      $S_{\neg}' = S_{\neg} \cap (\varphi_j|_{AP_j})$\;
      \If{$S' \neq \emptyset$ or $S_{\neg}' \neq \emptyset$}{
        add $(\pi,\, S',\, S_{\neg}')$ to $\textbf{Reply}_{i,j}^t$\;
      }
    }
  }
}
\textbf{return} $\textbf{Reply}_{i,j}^t$ for all reachable subscribers $j$\;
\end{algorithm}

Algorithm~\ref{ch03:alg:transferAlgo} summarizes the protocol. Requests are
subscription-based, while replies are event-driven and carry only new
subscriber-relevant information. The mechanism can be implemented over
multicast or unicast communication.

\subsubsection{Online plan verification and adaptation}
\label{ch03:subsec:multi-revise}

After receiving $\textbf{Sense}_i^t$ and replies
$\textbf{Reply}_{j,i}^t$ from neighbors, robot~$i$ updates
$\mathcal{T}_i^t$ as in Section~\ref{ch03:sec:partial}. Local sensing may
modify labels and adjacency information, whereas communicated replies modify
only labels, since they provide proposition-level knowledge rather than
traversability or cost information.
Plan verification and revision also follow Section~\ref{ch03:sec:partial}.
The sets $\widetilde{\Pi}_i^t$, $\Xi_i^t$, and $\aleph_i^t$ record updated
labels, invalid transitions, and unsafe transitions, respectively. If
$\Xi_i^t$ or $\aleph_i^t$ is nonempty, Algorithm~\ref{ch03:alg:revise} repairs
the run. Since local repair need not restore a cost-minimizing safe run, exact
synthesis is triggered periodically or after sufficient model change, in the
spirit of dynamic shortest-path updates
\cite{chan2009shortest,misra2005dynamic}. Let $\Upsilon_i^t$ count accumulated
model changes and $T_i$ record the last call to Algorithm~\ref{ch02:alg:optimal}.
When $\Upsilon_i^t \geq N_i^{\textup{call}}$ or
$t - T_i \geq T_i^{\textup{call}}$, the optimal search is called on
$\tilde{\mathcal{A}}_{r,i}^t$ from the compatible product-state set
$Q'_{i,\tau_{\textup{past}}}$, after which $\Upsilon_i^t$ and $T_i$ are reset.

\subsubsection{Overall structure}\label{ch03:subsec:structure}

Algorithm~\ref{ch03:alg:scheme} summarizes the collaborative planning
architecture. Each robot synthesizes an initial safe accepting run, exchanges
subscriptions and task-relevant sensing information during execution, updates
its transition system, and decides whether local revision or full optimal
synthesis is required. The next goal region is updated whenever the accepting
run changes or the current goal is reached.

\begin{algorithm}[t]
\caption{Cooperative online planning for each robot $i \in \mathcal{N}$}
\label{ch03:alg:scheme}
\DontPrintSemicolon
\KwIn{$\mathcal{T}_i^0$, $\tilde{\mathcal{A}}_{\varphi_i}$,
      $\tilde{\mathcal{A}}_{r,i}^{0}$, $x_i(t)$}
\KwOut{$R_i^t$, $\tau_i^t$, $\Upsilon_i^t$, $T_i$}

$R_i^0 = \texttt{OptRun}(\tilde{\mathcal{A}}_{r,i}^{0})$\;
$q'_{i,\textup{cur}} = q'_{i,\textup{next}} = R_{i,[\textup{pre},1]}^0$\;
$\pi_{i,\textup{next}} = q'_{i,\textup{next}}|_{\Pi_i}$\;
$R_{i,\textup{past}} = [\,]$, $\tau_{i,\textup{past}} = [\,]$\;
$\Upsilon_i^0 = 0$, $T_i = 0$\;

\While{\textup{\texttt{True}}}{
  send $\textbf{Request}_{i,g}^t$ to newly encountered neighbors
  $g \in \mathcal{N}_i^t$\;

  read $\textbf{Sense}_i^t$, incoming
  $\textbf{Request}_{j,i}^{t_0}$, and incoming
  $\textbf{Reply}_{h,i}^t$\;

  $\textbf{Reply}_{i,j}^t =
  \texttt{TranKnow}(\textbf{Sense}_i^t,\, \mathcal{S}_i^t)$\;

  send $\textbf{Reply}_{i,j}^t$ to reachable subscribers $j$\;

  $(\mathcal{T}_i^{t^+},\, \widetilde{\Pi}_i^t,\, \widehat{\Pi}_i^t) =
  \texttt{UpdateT}(\mathcal{T}_i^t,\,
  \textbf{Sense}_i^t,\, \textbf{Reply}_{h,i}^t)$\;

  $\tilde{\mathcal{A}}_{r,i}^{t^+} =
  \texttt{AdjProd}(\mathcal{T}_i^{t^+},\, \widehat{\Pi}_i^t,\,
  \tilde{\mathcal{A}}_{\varphi_i})$\;

  $(\aleph_i^t,\, \Xi_i^t) =
  \texttt{ValidRun}(\tilde{\mathcal{A}}_{r,i}^{t^+},\, R_i^t,\,
  \widetilde{\Pi}_i^t)$\;

  $\Upsilon_i^t = \Upsilon_i^t + |\widehat{\Pi}_i^t|$\;

  $Q'_{i,\tau_{\textup{past}}} =
  \texttt{CorProd}(\tilde{\mathcal{A}}_{r,i}^{t^+},\,
  \tau_{i,\textup{past}})$\;

  \If{$\Upsilon_i^t \geq N_i^{\textup{call}}$ or
      $t - T_i \geq T_i^{\textup{call}}$}{
    $R_i^{t^+} =
    \texttt{OptRun}(\tilde{\mathcal{A}}_{r,i}^{t^+},\,
    Q'_{i,\tau_{\textup{past}}})$\;
    $\Upsilon_i^t = 0$, $T_i = t$\;
  }
  \ElseIf{$\aleph_i^t \neq \emptyset$ or $\Xi_i^t \neq \emptyset$}{
    $R_i^{t^+} =
    \texttt{Revise}(\Xi_i^t,\, \aleph_i^t,\, R_i^t,\,
    \tilde{\mathcal{A}}_{r,i}^{t^+},\, Q'_{i,\tau_{\textup{past}}})$\;
  }
  \Else{
    $R_i^{t^+} = R_i^t$\;
  }

  \If{$x_i(t) \in q'_{i,\textup{next}}|_{\Pi_i}$ is confirmed}{
    $q'_{i,\textup{cur}} = q'_{i,\textup{next}}$\;
    $\tau_{i,\textup{past}} =
    \tau_{i,\textup{past}} + q'_{i,\textup{next}}|_{\Pi_i}$\;
    $R_{i,\textup{past}} =
    R_{i,\textup{past}} + q'_{i,\textup{next}}$\;
    $q'_{i,\textup{next}} =
    \texttt{NextGoal}(q'_{i,\textup{cur}},\, R_i^{t^+})$\;
  }

  $R_i^t = R_i^{t^+}$, $\tau_i^t = R_i^t|_{\Pi_i}$\;
  $\pi_{i,\textup{cur}} = q'_{i,\textup{cur}}|_{\Pi_i}$\;
  $\pi_{i,\textup{next}} = q'_{i,\textup{next}}|_{\Pi_i}$\;
  $u_i = \mathcal{U}(x_i(t),\, \pi_{i,\textup{cur}},\,
  \pi_{i,\textup{next}})$\;
}
\end{algorithm}

\begin{theorem}\label{ch03:theo:improvement}
For each robot $i \in \mathcal{N}$ and each time $t \geq 0$, the plan
$\tau_i^t$ produced by Algorithm~\ref{ch03:alg:scheme} is valid and safe with
respect to the current transition system $\mathcal{T}_i^t$ and the local hard
specification $\varphi_i^{\textup{hard}}$. Moreover, for any $t' \geq 0$,
there exists some $t \in [t',\, t' + T_i^{\textup{call}}]$ such that
$\tau_i^t$ is induced by a safe accepting run of
$\tilde{\mathcal{A}}_{r,i}^t$.
\end{theorem}

\begin{proof}
The first part follows from the update and revision mechanisms established in
Section~\ref{ch03:sec:partial}. If no invalid or unsafe transition is detected,
the current plan is retained and remains valid and safe by
Theorem~\ref{ch03:theo:validity}. If invalid or unsafe transitions are
detected, Algorithm~\ref{ch03:alg:revise} returns a compatible continuation,
and Theorem~\ref{ch03:theo:new-safe} guarantees validity and safety of the
resulting complete trajectory. Whenever
Algorithm~\ref{ch02:alg:optimal} is called, it computes a safe accepting run of
$\tilde{\mathcal{A}}_{r,i}^t$ from the compatible initial-state set
$Q'_{i,\tau_{\textup{past}}}$, so the induced plan is again valid and safe.
For the second part, the time-trigger condition
$t - T_i \geq T_i^{\textup{call}}$ guarantees that
Algorithm~\ref{ch02:alg:optimal} is invoked at least once within any interval
of length $T_i^{\textup{call}}$. At each such invocation, the resulting plan is
induced by a safe accepting run of the current product automaton. This proves
the claimed recurrence property.
\end{proof}
